\documentclass[10pt,a4paper]{amsart}
\usepackage[margin=19mm]{geometry}
\usepackage{amsmath,amssymb,mathtools}
\usepackage[scr=rsfs,cal=euler]{mathalfa}
\usepackage{xfrac}
\usepackage[unicode,hidelinks,bookmarks=false]{hyperref}
\newcommand{\ie}{i.e.\ }
\newcommand{\cf}{cf.\ }
\newcommand{\resp}{respectively\ }
\newcommand{\Subsection}{\subsection}
\providecommand{\nobreakdash}{\nobreak\hbox{-}}
\setlength{\emergencystretch}{2em}
\multlinegap0pt
\usepackage{bm}
\usepackage{bbm}
\usepackage{dsfont}
\usepackage{xspace}
\usepackage{cancel}
\usepackage{mathtools}
\usepackage{cleveref}
\usepackage{amsthm}
\makeatletter
\@ifundefined{theorem}{\newtheorem{theorem}{Theorem}}{}
\makeatother
\makeatletter
\newcommand{\Ac}{\mathcal{A}}
\newcommand{\Bc}{\mathcal{B}}
\newcommand{\Bf}{\mathfrak{B}}
\newcommand{\Hf}{\mathfrak{H}}
\newcommand{\Zf}{\mathfrak{Z}}
\@ifundefined{boxedtikzcd}{\newenvironment{boxedtikzcd}
 {\begin{lrbox}{\boxedtikzcdbox}\begin{tikzcd}}}{}
\makeatother
\title[Approximation properties of double complexes]{Approximation properties of double complexes}
\author{Daniel F{\o}rland Holmen, Jan Martin Nordbotten, Jon Eivind Vatne}
\date{Source: arXiv:2604.13982v1 (CC BY 4.0)}
\begin{document}
\maketitle

\noindent\emph{Source and licence.} Approximation properties of double complexes. Version arXiv:2604.13982v1 (CC BY 4.0), distributed under the Creative Commons Attribution 4.0 International licence, as recorded in the arXiv licence metadata for this exact version and shown on the abstract page https://arxiv.org/abs/2604.13982v1. Full rights notice: licenses/arxiv-2604.13982-CC-BY-4.0.txt.\par\smallskip

\noindent\textit{Editorial scope.} Counted unit: the Theorem whose source label is \texttt{theorem: existence of bcp} in the article (source0.tex lines 466--471, source label \texttt{theorem: existence of bcp}), delivered together with the article's own complete original proof of it (source0.tex lines 473--526, one \texttt{proof} environment of 54 lines). No mathematics is written, repaired or re-derived: statement and proof are the article's own text line for line. Required context retained uncounted: eq: assumpions (lines 457--462); eq: ass1 (lines 458--461); eq: ass2 (lines 458--461). Every definition, display and label of those lines is retained unchanged. Omitted from this package and not redistributed: 1--456, 462--465, 472--472, 527--1355. Nothing inside a retained excerpt is omitted or altered: every retained body is delivered exactly as the article's lines are written, so no omission receipt of any kind accompanies this package. Citations: the retained excerpts carry no citation command at all, so no bibliography entry of the article is transcribed and no reference is redistributed. Reference adaptation: none was needed and none was made; every reference inside the delivered unit resolves inside it, so no unresolved reference or citation is produced. Printed slips kept literal: no printed word, symbol, hypothesis or number of the counted statement or of its proof was altered. Layout and class: the article's own class is \texttt{amsart} with packages including caption, subcaption, booktabs, siunitx, inputenc, hyperref, graphicx, bbm, soul, tikz, mathrsfs, amsmath, amsthm, tikz-cd; the wrapper is the lane's pinned \texttt{amsart} editorial wrapper at 10pt on A4, which restates the theorem-like environments the retained text needs and adds the article's own macro definitions that the retained text needs (\texttt{\textbackslash Ac}, \texttt{\textbackslash Bc}, \texttt{\textbackslash Bf}, \texttt{\textbackslash Hf}, \texttt{\textbackslash Zf}), with extra wrapper packages (bm, bbm, dsfont, xspace, cancel, mathtools, cleveref). The delivered numbering is the unit's own. Assets: no retained span contains a figure environment, a graphics-inclusion command, a PDF-page inclusion, a table or a TikZ picture; no image file is copied into this package, the package declares no asset dependency and no image file is redistributed with it. Licence: version arXiv:2604.13982v1 of this article is distributed under the Creative Commons Attribution 4.0 International licence, as recorded in the arXiv licence metadata for this exact version and shown on the abstract page https://arxiv.org/abs/2604.13982v1. A full rights notice is delivered at licenses/arxiv-2604.13982-CC-BY-4.0.txt. Characters: every retained excerpt is pure ASCII, so the delivered engine, XeLaTeX with its default Latin Modern text font, reports no missing character.
\begin{subequations}\label{eq: assumpions}
\begin{align}\label{eq: ass1}
\|v\|_\Ac &\leq \kappa_1 \|dv\|_\Ac, \qquad &v \in \mathfrak{Z}^{k,\perp}(\Bc), \\
\|q\|_\Ac &\leq \kappa_2 \|P_{\mathfrak{H}(\Ac)} q\|_\Ac, \qquad &q \in \Hf^k(\Bc). \label{eq: ass2}
\end{align}   
\end{subequations}

\begin{theorem}\label{theorem: existence of bcp}
Suppose the assumptions in \cref{eq: assumpions} hold. Then there exists bounded cochain projections $\pi^k: \Ac^k \to \Bc^k$, where $\|\pi^k\|$ is bounded by the constants $\kappa_1$ and $\kappa_2$:
\begin{equation}
\|\pi^k_\Bc v \|_\Ac \leq (1 + \kappa_1 + \kappa_1 \kappa_2) \|v\|_\Ac, \qquad \forall v \in \Ac^k.
\end{equation}
\end{theorem}

\begin{proof}
The first assumption implies that there exists an operator
\begin{equation}
Q_\Bc: \Ac^k \to \mathfrak{Z}^{k, \perp}(\Bc), \qquad dQ_\Bc v = P_{\mathfrak{B}(\Bc)} dv.
\end{equation}
It follows from \cref{eq: ass1} that $Q_\Bc$ is bounded:
\begin{equation}
\|Q_\Bc v\|_\Ac \leq \kappa_1 \|d Q_\Bc v\|  = \kappa_1\| P_{\mathfrak{B}(\Bc)} dv \| \leq \kappa_1
\|v\|_\Ac.
\end{equation}

Next, we consider the orthogonal projection $P_\Hf: \Ac^k \to \Hf^k$. When restricting the domain of the projection to $\Hf(\Bc)$, it defines a map $P_\Hf|_{\Hf(\Bc)}: \Hf^k(\Bc) \to \Hf^k$. We denote the inverse of this map by $R_\Bc: \Hf \to \Hf(\Bc)$, which exists and is bounded because of \cref{eq: ass2}.

We now consider the following map:
\begin{equation}\label{eq: bounded cochain projection Q and R}
\pi_\Bc^k = P_{\mathfrak{B}(\Bc)} + R_\Bc P_{\mathfrak{H}}(I - Q_\Bc) + Q_\Bc. 
\end{equation}
We claim that \cref{eq: bounded cochain projection Q and R} defines a bounded cochain projection. For this to be true, we need to verify the following three properties for \cref{eq: bounded cochain projection Q and R}: 
\begin{enumerate}
    \item $\pi_\Bc^k$ is a cochain map
    \item $\pi_\Bc^k$ is bounded
    \item $\pi_\Bc^k$ is a projection
\end{enumerate}
Firstly, for $\pi^k_\Bc$ to be a cochain map it satisfies $d\pi^k_\Bc v = \pi^{k+1}_\Bc dv$. Using the Hodge decomposition, we can decompose $v = du + q + d^*w$.

On one hand, 
\begin{equation}
d\pi_\Bc^{k+1} v = dP_{\Bf_\Bc} v + d(R_\Bc P_\Hf (I - Q_\Bc ) v) + dQ_\Bc v.
\end{equation}
The first two terms are zero, since $P_{\Bf_\Bc} v = du_\Bc$ and thus $dP_{\Bf_\Bc} v= ddu_\Bc = 0$, and similarly $dq_\Bc = 0$.

On the other hand, 
\begin{equation}
\pi_\Bc^k dv = P_{\Bf_\Bc} dv + (R_\Bc P_\Hf (I - Q_\Bc) dv) + Q_\Bc dv.
\end{equation}
In this expression, the second term and the third term are zero. The second term is zero since $P_\Hf dv = 0$, while the last term is zero since $Q_\Bc: \Ac^k \to \Zf^{k,\perp}(\Bc)$, but $dv \in \Bf^k \subset \Zf^k$, hence $Q_h dv = 0$. By definition of $Q_\Bc$, $dQ_\Bc v = P_{\Bf_\Bc} dv$. The only non-zero terms coincide, and thus we conclude that $\pi_\Bc^{k+1} dv = d\pi^k_\Bc v$.


Secondly, $\pi^k_\Bc$ is bounded because both $Q_\Bc$ and $R_\Bc$ are bounded, and each orthogonal projection $P_K$ is bounded with bound 1.

\begin{subequations}
\begin{align}
\| \pi_\Bc^k v\|_\Ac &= \|(P_{\mathfrak{B}(\Bc)} + R_\Bc P_{\mathfrak{H}}(I - Q_\Bc) + Q_\Bc)v \|_\Ac \\
&\leq  \| P_{\mathfrak{B}(\Bc)}v \|_\Ac + \| R_\Bc P_{\mathfrak{H}}(I - Q_\Bc)v \|_\Ac + \| Q_\Bc v \|_\Ac \\ 
&\leq \|v\|_\Ac + \kappa_1 \kappa_2 \|v\|_\Ac + \kappa_1 \|v\|_\Ac = (1 + \kappa_1 + \kappa_1 \kappa_2) \|v\|_\Ac.
\end{align}
\end{subequations}

If $v_\Bc = du_\Bc + q_\Bc + d^*_\Bc w_\Bc \in \Bc^k$, then $Q_\Bc v_\Bc = P_{\Bf^*(\Bc)} v_\Bc$ and $R_\Bc P_{\Hf} v_\Bc = P_{\Hf(\Bc)} v_\Bc$. From this we see that the map $\pi_\Bc^k$ is invariant on $v_\Bc \in \Bc^k$: 
\begin{equation}
\pi_\Bc^k v_\Bc = du_\Bc + q_\Bc + d^*_\Bc w_\Bc = v_\Bc.
\end{equation}
This invariance proves the final claim that $\pi^k$ is a projection. We can therefore conclude that $\pi^k$ as defined in \cref{eq: bounded cochain projection Q and R} is a bounded cochain projection.
\end{proof}

\end{document}
